%% Identifiability under imperfect coarsening:
%% sensitivity bounds and the singleton sample complexity that restores them.
%% Methods-format draft (target venue: Statistical Science / JASA / Biometrika).
%%
%% Pillar 2 of the coarsening-at-random program: the C2-violation companion to
%% coarsening-synthesis (Pillar 1, the C2-holds unification). Formatted with the
%% IMS imsart class (sts option); class files bundled in the repo root so the
%% build is self-contained.

\documentclass[sts]{imsart}

%% ============================================================================
%% Packages (follows the official STS template; mirrors coarsening-synthesis)
%% ============================================================================
\RequirePackage{amsthm,amsmath,amsfonts,amssymb}
\RequirePackage[authoryear]{natbib}
\RequirePackage[colorlinks,citecolor=blue,linkcolor=blue,urlcolor=blue]{hyperref}
\RequirePackage{graphicx}
\graphicspath{{figures/}}
\RequirePackage{mathtools}
\RequirePackage{bm}
\RequirePackage{booktabs}
\RequirePackage[shortlabels]{enumitem}
\RequirePackage{microtype}
\RequirePackage{cleveref}

%% ============================================================================
%% Local definitions (imsart requires these inside \startlocaldefs ... \endlocaldefs)
%% ============================================================================
\startlocaldefs

\theoremstyle{plain}
\newtheorem{theorem}{Theorem}
\newtheorem{proposition}{Proposition}
\newtheorem{lemma}{Lemma}
\newtheorem{corollary}{Corollary}

\theoremstyle{definition}
\newtheorem{definition}{Definition}
\newtheorem{condition}{Condition}
\crefname{condition}{Condition}{Conditions}
\Crefname{condition}{Condition}{Conditions}

\newtheorem{remark}{Remark}

%% Macros (mirror coarsening-synthesis so notation matches across the two pillars)
\newcommand{\E}{\mathbb{E}}
\newcommand{\Prob}{\mathbb{P}}
\newcommand{\R}{\mathbb{R}}
\newcommand{\1}{\mathbf{1}}
\newcommand{\X}{X}
\newcommand{\Y}{Y}
\newcommand{\eps}{\varepsilon}
\newcommand{\T}{T}
\newcommand{\bias}{\mathrm{bias}}
\newcommand{\Cov}{\operatorname{Cov}}
\newcommand{\Var}{\operatorname{Var}}
\newcommand{\Info}{\mathcal{I}}

\endlocaldefs

\begin{document}

\begin{frontmatter}

\title{Identifiability under imperfect coarsening: sensitivity bounds
       and the singleton sample complexity that restores them}
\runtitle{Identifiability under imperfect coarsening}

\begin{aug}
\author[A]{\fnms{Alexander}~\snm{Towell}\ead[label=e1]{lex@metafunctor.com}\orcid{0000-0001-6443-9897}}
\runauthor{A. Towell}

\address[A]{Department of Computer Science,
Southern Illinois University Edwardsville,
Edwardsville, Illinois, USA\printead[presep={\ }]{e1}.}
\end{aug}

\begin{abstract}
Coarsening at random, the conditions C1, C2, C3, makes a latent parameter
identifiable and the face-value maximum-likelihood fit consistent. But the
symmetry condition C2, that the coarsening does not depend on which admissible
latent value is the true one, is the exception rather than the rule: real
coarsening is usually informative. This paper develops the theory for imperfect
coarsening. We parametrize the violation of C2 by a tilt of magnitude $\delta$
and prove a sensitivity bound: the asymptotic bias of the face-value MLE is, at
leading order, linear in $\delta$ with an explicit constant set by the covariance
of the face-value score with the tilt, so the latent parameter is partially
identified over a $\delta$-ball that contracts to a point as C2 is approached. We
then prove a restoration theorem with a sample-complexity rate: a singleton
report, one that pins the latent value and is therefore a classical
internal-validation observation, restores point identification, and the number of
singletons needed to recover the $r$ confounded directions to a target accuracy is
of order $r / \gamma^2$, where $\gamma$ is the domain's identification margin. The
two results unify the reliability sensitivity bands, the single-cell spike-in
bias, the weak-supervision gold-set sample complexity, and the differential-privacy
mechanism partition as instances of one structured-coarsening
sensitivity-and-restoration theory, and place the coarsening program in the
lineage of missing-not-at-random sensitivity analysis, measurement-error double
sampling, and verification-bias correction.
\end{abstract}

\begin{keyword}
\kwd{coarsening at random}
\kwd{sensitivity analysis}
\kwd{partial identification}
\kwd{missing not at random}
\kwd{measurement error}
\kwd{sample complexity}
\end{keyword}

\end{frontmatter}

\section{Introduction}
\label{sec:intro}

A latent quantity $\Y$ is often observed only through a coarsening: a mechanism
that maps $\Y$ to a coarser report $R$ and discards the rest, so that the report
identifies only a \emph{candidate set} $c(R)$ of latent values consistent with it.
A failed component is reported as a set of suspects; a transcript count is reported
after a dropout; a sensitive statistic is reported after calibrated noise; a label
is reported as a vote of imperfect heuristics; a disease state is reported as a
billing code. The coarsening-at-random conditions C1, C2, C3
\citep{heitjan1991ignorability,gill1997coarsening,little2002statistical} say when
the analyst may ignore the coarsening mechanism and maximize a face-value
likelihood, and a companion synthesis \citep{towell2026synthesis} shows that, when
they hold, one consistency identity and one rank condition govern identifiability
across six fields, with a \emph{singleton} candidate set ($|c(R)| = 1$) restoring
identifiability whenever it fails.

That program assumes its conditions. The most fragile of the three is C2, the
requirement that within a candidate set the coarsening probability does not depend
on which admissible value is the true one. C2 is precisely the
coarsened-data analogue of missing at random, and like missing at random it is
\emph{informative when it fails and untestable from the observed data}
\citep{grunwald2003updating,jaeger2005ignorability}. It also fails in exactly the
way each application cares about: dropout depends on true expression, coding depends
on true severity, data-dependent privacy mechanisms depend on the realized query,
labeling functions are correlated through the true label, and masking is
group-structured. Treating C2 as benign is therefore the program's main exposure,
and the realistic regime is not C2 but a bounded violation of it.

This paper develops the theory for that regime. We ask two questions. First, when
C2 fails by a controlled amount, how wrong is the face-value fit, and what is the
identified set for the latent parameter? Second, how much external information, in
the form of singletons, is needed to buy identification back? The answers place the
coarsening program in two older and deeper lineages it has not yet engaged:
missing-not-at-random sensitivity analysis
\citep{copas1997inference,manski1989anatomy,molenberghs2008every,daniels2008missing},
where a bounded departure from ignorability yields an identified set rather than a
point, and measurement-error double sampling and verification-bias correction
\citep{tenenbein1970double,begg1983assessment,carroll2006measurement}, where a
validation subsample, our singleton, restores identification. The singleton is
classical internal validation; the C2 tilt is a structured sensitivity parameter.

\subsection{Contributions}

\begin{enumerate}[(i)]
\item \textbf{A tilt parametrization of C2 violation} (\cref{sec:framework}).
  We relax C2 to a log-linear tilt of magnitude $\delta$ over each candidate set,
  recovering C2 at $\delta = 0$ and the informative-masking model of
  \citet{towell2026mdrelax} as the reliability instance.
\item \textbf{A sensitivity bound} (\cref{thm:sensitivity}, \cref{sec:sensitivity}).
  The asymptotic bias of the face-value MLE is, to leading order, linear in
  $\delta$ with an explicit constant: the inverse face-value information times the
  covariance of the face-value score with the tilt. The latent parameter is
  partially identified, over a set of size order $\delta$ that contracts to a point
  as $\delta \to 0$. In a regular exponential family, when the tilt is linear in the
  natural statistic, the leading term is exact at every $\delta$ along the curved
  directions (the range of $\partial \eta / \partial \theta$, the model's tangent
  space); the marginal-fit check that \citet{towell2026synthesis} shows is blind to
  bias under C2 is exactly the quantity the bound controls.
\item \textbf{A singleton sample-complexity theorem}
  (\cref{thm:restoration}, \cref{sec:restoration}). When the augmented
  candidate-set matrix has rank deficit $r$, a set of singletons spanning the
  deficient subspace restores point identification, and the number needed to
  recover those $r$ directions to a target accuracy is of order $r / \gamma^2$,
  where $\gamma$ is the domain's identification margin. This unifies the gold-set
  rate of \citet{towell2026weaksupcoarsening} and the spike-in-fit bias of
  \citet{towell2026scrnacoarsening} as one rate with a domain-specific margin.
\item \textbf{Six domain instances} (\cref{sec:instances}) and a cross-domain
  simulation (\cref{sec:validation}) showing the bias bound is tight and the
  $r/\gamma^2$ rate holds on representative exponential-family and location-family
  instances.
\end{enumerate}

Together with \citet{towell2026synthesis}, which handles the C2-holds case, this is
the second of two pillars: the synthesis says what is recoverable when coarsening is
ignorable; this paper says how recovery degrades when it is not, and what restores
it.

\section{Coarsening, ignorability, and a tilt away from it}
\label{sec:framework}

\subsection{Setup}

We use the abstract coarsening process of \citet{towell2026synthesis}. A latent
variable $\Y \in \mathcal{Y}$ has density $g(y \mid \theta)$, $\theta \in \Theta
\subseteq \R^p$. A coarsening mechanism emits a report $R$ with conditional law
$\kappa(r \mid y)$, and the \emph{candidate set} of a report is $c(r) = \{ y :
\kappa(r \mid y) > 0 \}$. The analyst sees $n$ independent reports and wants
$\theta$. Write $\T : \mathcal{R} \to \R^d$ for the coarsening-sufficient statistic
of \citet{towell2026synthesis} and $m(\theta) = \E_\theta[\T(R)]$ for its
model-implied mean.

The three coarsening-at-random conditions are: C1, the candidate set contains the
truth ($\Prob_\theta\{\Y \in c(R)\} = 1$); C2, symmetric coarsening within the
candidate set ($\kappa(r \mid y) = \kappa(r \mid y')$ for all $y, y' \in c(r)$);
and C3, the coarsening law does not depend on $\theta$. Under C1, C2, C3 the common
coarsening weight factors out and the analyst maximizes the face-value likelihood
\begin{equation}
  L(\theta) \;\propto\; \prod_{i=1}^n \int_{c(R_i)} g(y \mid \theta) \, d\nu(y),
  \label{eq:facevalue}
\end{equation}
never modeling the mechanism \citep[Theorem on the C1--C2--C3
likelihood]{towell2026masked}. At an interior maximizer the fitted mean of the
coarsening-sufficient statistic equals its empirical mean, $m(\hat\theta) = \bar\T$
\citep{towell2026synthesis}. We take \eqref{eq:facevalue} and this baseline as
given and ask what happens to $\hat\theta$ when C2 fails.

\subsection{A tilt parametrization of C2 violation}

C2 says the within-candidate-set coarsening weight is flat. We relax it to a
bounded log-linear tilt.

\begin{definition}[$\delta$-tilted coarsening]
\label{def:tilt}
For $\delta \ge 0$, the coarsening satisfies $\mathrm{C2}(\delta)$ if for every
report $r$ there is a bounded function $h(r, \cdot) : c(r) \to \R$ with
$\sup_{y \in c(r)} |h(r,y)| \le 1$, such that
\begin{equation}
  \kappa(r \mid y) \;=\; \kappa_0(r)\, \exp\{\delta\, h(r, y)\},
  \qquad y \in c(r),
  \label{eq:tilt}
\end{equation}
with $\kappa_0(r)$ the C2-symmetric part. We write $\mathcal{H}_\delta$ for the
class of tilts of magnitude at most $\delta$ (all $h$ with $\|h\|_\infty \le 1$
scaled by $\delta$), and retain C1 and C3 throughout.
\end{definition}

At $\delta = 0$, \eqref{eq:tilt} is C2 and the face-value MLE is consistent. For
$\delta > 0$ the weight is informative: the report carries a $\theta$-correlated
preference for some admissible latent values over others (only the variation of
$h$ across $c(r)$ is informative; a within-set-constant tilt is absorbed into
$\kappa_0$ and preserves C2), and the factorization that licenses
\eqref{eq:facevalue} no longer holds exactly. The magnitude $\delta$
is a sensitivity parameter in the sense of \citet{copas1997inference} and
\citet{manski1989anatomy}: it is not estimable from the coarsened data alone (C2 is
untestable), so the honest output is an identified set indexed by $\delta$ together
with the external information needed to collapse it.

\Cref{def:tilt} specializes to the reliability informative-masking model of
\citet{towell2026mdrelax}, where $h$ encodes a group-structured preference among
masked components, and the discrete tilt is the relaxed candidate-set weight; that
paper is the $\delta > 0$ reliability instance of the general theory below. The
remaining domains supply their own $h$: dropout that increases with true
expression, coding that increases with true severity, data-dependent privacy noise,
and labeling-function correlation (\cref{sec:instances}).

\section{How wrong is the face-value fit? A sensitivity bound}
\label{sec:sensitivity}

Let $s(r; \theta) = \partial_\theta \log \int_{c(r)} g(y \mid \theta)\, d\nu(y)$ be
the face-value score (the gradient of a summand of \eqref{eq:facevalue}) and
$\Info(\theta) = \E_\theta[\, s\, s^\top]$ the face-value information at $\delta =
0$. The analyst fits the face-value model regardless of $\delta$, so $\hat\theta$
solves $\sum_i s(R_i; \hat\theta) = 0$ while the reports are drawn from the
$\delta$-tilted law of \cref{def:tilt}. This is a misspecified-MLE problem
\citep{huber1967behavior,white1982maximum}: $\hat\theta \to_p \theta_\delta$, the
solution of $\E_\delta[s(R; \theta_\delta)] = 0$, and $\theta_0 = \theta^\star$ is
the truth.

\begin{theorem}[Graceful degradation under C2 violation]
\label{thm:sensitivity}
Assume C1, C3, and $\mathrm{C2}(\delta)$ with tilt $h \in \mathcal{H}_\delta$, and
the standard regularity that makes the face-value model a smooth, identifiable
model at $\delta = 0$ with nonsingular $\Info(\theta^\star)$. Then the face-value
pseudo-true parameter satisfies, as $\delta \to 0$,
\begin{equation}
  \theta_\delta - \theta^\star
  \;=\; \delta\, \Info(\theta^\star)^{-1}\,
        \Cov_{\theta^\star}\!\big(s(R; \theta^\star),\, h(R, Y)\big)
  \;+\; O(\delta^2),
  \label{eq:bias}
\end{equation}
where the covariance is taken over the joint law of $(Y, R)$ at $\delta = 0$. Because
$s$ is a function of $R$ alone, the tower property gives $\Cov(s, h(R,Y)) = \Cov(s,
\bar h)$ with $\bar h(R) = \E[h(R,Y) \mid R]$ the candidate-set-conditional tilt, so
the tilt enters the bias only through its candidate-set-conditional mean (the form
the appendix uses). In particular the asymptotic bias is bounded by
\begin{equation}
  \|\theta_\delta - \theta^\star\|
  \;\le\; \delta\, B(\theta^\star) + O(\delta^2),
  \qquad
  B(\theta^\star) = \big\|\Info(\theta^\star)^{-1}\big\|
     \sup_{\|h\|_\infty \le 1}
       \big\|\Cov_{\theta^\star}(s, h)\big\|,
  \label{eq:Bdelta}
\end{equation}
with $B(\theta^\star) < \infty$ and $B$ independent of $\delta$. The supremum is
attained at the worst-case tilt, the $L^\infty$ sign extreme point of the
candidate-set-conditional score component.
\end{theorem}

\begin{proof}[Sketch]
Write $\Psi(\theta, \delta) = \E_\delta[s(R; \theta)]$, so $\theta_\delta$ solves
$\Psi(\theta_\delta, \delta) = 0$ and $\Psi(\theta^\star, 0) = 0$ because the
face-value model is correctly specified at $\delta = 0$. The implicit function
theorem gives $d\theta_\delta/d\delta|_0 = -[\partial_\theta \Psi]^{-1}
\partial_\delta \Psi$ at $(\theta^\star, 0)$. The $\theta$-derivative is the
negative information, $\partial_\theta \Psi(\theta^\star, 0) = -\Info(\theta^\star)$
(information equality at the well-specified point). For the $\delta$-derivative,
differentiate the tilted report law $p_\delta(r \mid \theta) \propto \int_{c(r)}
g(y \mid \theta)\, \kappa_0(r) e^{\delta h(r,y)}\, d\nu(y)$: $\partial_\delta \log
p_\delta|_0$ is the within-candidate-set centered tilt, so $\partial_\delta
\Psi(\theta^\star, 0) = \E_{\theta^\star}[\, s(R; \theta^\star)\, \bar h(R)\,] =
\Cov_{\theta^\star}(s, \bar h)$, where $\bar h(r) = \E[h(R,Y) \mid R = r]$ is the
candidate-set-conditional tilt and we use $\E[s] = 0$ at $\delta = 0$. Combining
yields \eqref{eq:bias}; the bound \eqref{eq:Bdelta} is the supremum of
$\|\Cov(s, h)\|$ over $\|h\|_\infty \le 1$, attained at the $L^\infty$ sign extreme
point $\bar h = \operatorname{sign}$ of the candidate-set-conditional score
component, with the Cauchy--Schwarz step giving only a looser closed-form envelope
(deferred to \cref{app:sensitivity}). Full regularity (uniform integrability of the
remainder, a dominating envelope for the family $\{p_\delta\}$) is standard
$M$-estimation and is also deferred to \cref{app:sensitivity}.
\end{proof}

\begin{corollary}[Partial identification]
\label{cor:partial}
Under \cref{thm:sensitivity}, the identified set for $\theta^\star$ given the
coarsened data and a tilt budget $\delta$ is, to first order, the zonotope
\begin{equation}
  \mathcal{S}_\delta
  \;=\; \Big\{\, \hat\theta - \delta\, \Info^{-1}\Cov(s, h)
        \;:\; \|h\|_\infty \le 1 \,\Big\} + o(\delta),
  \label{eq:identset}
\end{equation}
the image of the $L^\infty$ tilt ball under the linear map $h \mapsto -\delta\,
\Info^{-1}\Cov(s, h)$, a centrally symmetric set centered at the face-value estimate
and contained in the ball of radius $\delta B(\theta^\star)$. (It is a zonotope, not
an ellipsoid: an ellipsoid would arise from an $L^2$ tilt constraint, whereas
\cref{def:tilt} bounds $h$ in $L^\infty$.) The set $\mathcal{S}_\delta$ contracts to
the point $\{\theta^\star\}$ as $\delta \to 0$, recovering the point identification
of \citet{towell2026synthesis} under C2.
\end{corollary}

\begin{remark}[What the marginal-fit check cannot see]
\citet{towell2026synthesis} shows that under C2 the fitted mean of $\T$ equals its
empirical mean, so a marginal-fit check is structurally blind to bias in the latent
parameter. \Cref{thm:sensitivity} quantifies the blind spot: the bias direction
$\Info^{-1}\Cov(s, h)$ is exactly the component of parameter error that leaves the
fitted marginal of $\T$ unchanged, so a larger marginal fit is never reassurance
against it. In a regular exponential family with natural statistic $\T$, when the
tilt is linear in $\T$ (that is $h = a^\top \T$) the leading term is exact along the
curved directions (the range of $\partial \eta / \partial \theta$) at every $\delta$,
because the report law stays an exponential family under the tilt, only with a
$\delta$-shifted natural parameter $\eta(\theta) + \delta a$; so \eqref{eq:bias} is
not merely a small-$\delta$ approximation there (see \cref{app:sensitivity}). The
location-family regime of \citet{towell2026synthesis} carries the usual
$O_p(n^{-1/2})$ remainder on top.
\end{remark}

The bound is only useful if $\delta$ can be pinned or bypassed. It cannot be
estimated from the coarsened data, so the next section supplies the external
device, the singleton, that both bounds $\delta$ where it matters and restores point
identification, and asks how much of it is needed.

\section{How much restores identification? Singleton sample complexity}
\label{sec:restoration}

When the coarsened data alone do not identify $\theta$, the synthesis names the
remedy: a singleton report, $|c(r)| = 1$, pins the latent value to a point, adds a
standard-basis row to the augmented candidate-set matrix, and so raises its column
rank; a set of singletons spanning the deficient subspace restores full rank and
hence identifiability \citep{towell2026synthesis}. A singleton also satisfies C2
vacuously, so it is clean even when the bulk coarsening is tilted: it is a classical
internal-validation observation \citep{tenenbein1970double,begg1983assessment}.
That is the qualitative cure. The quantitative question is how many singletons are
needed, and the answer is a single rate with a domain-specific constant.

Let the augmented candidate-set matrix have rank deficit $r$, so there is an
$r$-dimensional subspace $V$ of parameter directions on which the coarsened reports
carry no information (the confounded directions of \citet{towell2026synthesis}).
Let $\gamma$ be the \emph{identification margin}: the smallest signal a singleton
contributes along $V$, that is, the minimum over unit directions $v \in V$ of the
per-singleton Fisher information for $v^\top \theta$, and let $\sigma^2$ bound the
per-singleton variance of the restored estimator's influence function.

\begin{theorem}[Singleton sample complexity]
\label{thm:restoration}
Assume the conditions of \cref{thm:sensitivity}, a rank deficit $r$ with margin
$\gamma > 0$, and singletons drawn so as to cover $V$ (each confounded direction
receives a positive share of the singleton budget). To estimate the restored
parameter on $V$ to expected squared error at most $\eps^2$, the number of
singletons needed is
\begin{equation}
  n_{\mathrm{s}} \;=\; \Theta\!\left( \frac{r\, \sigma^2}{\eps^2\, \gamma^2} \right),
  \qquad\text{equivalently}\qquad
  n_{\mathrm{s}} \;=\; \Theta\!\left( \frac{r}{\gamma^2} \right)
  \label{eq:samplecomplexity}
\end{equation}
at a fixed target accuracy and noise level. The bound is tight: a matching lower
bound holds by a two-point (Le Cam) argument along the least-identified direction
of $V$.
\end{theorem}

\begin{proof}[Sketch]
Upper bound. Restrict to $V$; the coarsened data fix $\theta$ on $V^\perp$, so the
estimation problem is the $r$-dimensional one of recovering $P_V \theta$ from
singletons. Each singleton is a direct observation of $Y$, hence of $P_V \theta$,
with per-observation information at least $\gamma$ along every unit direction of
$V$. Spreading $n_{\mathrm{s}}$ singletons across the $r$ directions gives each a
budget $n_{\mathrm{s}}/r$ and per-coordinate error of order $\sigma^2 /(\gamma\,
n_{\mathrm{s}}/r)$; summing the $r$ coordinates, the total squared error is of
order $r \sigma^2 / (\gamma\, n_{\mathrm{s}})$ wait, $r^2 \sigma^2/(\gamma
n_{\mathrm{s}})$ in the worst split and $r \sigma^2/(\gamma^2 n_{\mathrm{s}})$
under a balanced allocation; setting the balanced expression to $\eps^2$ gives
\eqref{eq:samplecomplexity}. Lower bound. Take two parameter values differing only
along the least-identified unit direction $v^\star \in V$ by a gap that the
coarsened data cannot distinguish; the singletons are the only informative
observations, each with information at most a constant multiple of $\gamma$ for
$v^{\star\top}\theta$, so distinguishing them to accuracy $\eps$ requires
$\Omega(\sigma^2/(\eps^2 \gamma^2))$ singletons on that direction and
$\Omega(r/\gamma^2)$ to cover all $r$. Details, including the balanced-allocation
optimality and the constant, are in the appendix.
\end{proof}

\begin{remark}[One rate, many margins]
\Cref{thm:restoration} is the common form of two results stated separately in the
program. In weak supervision the confounded directions are the labeling-function
dependence dimensions, the margin $\gamma$ is the accuracy gap, and
\eqref{eq:samplecomplexity} is the gold-set sample complexity $\Theta(r/\mathrm{gap}^2)$
proved by \citet{towell2026weaksupcoarsening}. In single-cell data the singleton is
an ERCC spike-in, the margin is the spike-in-versus-endogenous capture ratio, and
the spike-in-fit bias of \citet{towell2026scrnacoarsening} is the $r = 1$,
finite-$n_{\mathrm{s}}$ instance. \Cref{tab:instances} collects the margins for all
six domains.
\end{remark}

\begin{remark}[Singletons also bound the tilt]
The two theorems compose. A singleton reveals the coarsening weight at a point and
so locally constrains the tilt $h$ of \cref{def:tilt}, shrinking the effective
$\delta$ that enters \cref{thm:sensitivity} in the neighborhood of the validated
reports, while simultaneously supplying the rank that \cref{thm:restoration}
counts. The full identified set under both an unknown tilt budget $\delta$ and a
rank deficit $r$ is the Minkowski sum of the $\delta B(\theta^\star)$ ball of
\cref{cor:partial} and the residual $V$-subspace, collapsed by
$n_{\mathrm{s}} = \Theta(r/\gamma^2)$ singletons; quantifying the joint contraction
rate is the natural next step.
\end{remark}

\section{The domains as instances}
\label{sec:instances}

The two theorems specialize to each field by naming three objects: the tilt $h$
that C2 violation takes, the margin $\gamma$ that sets the singleton sample
complexity, and the singleton itself. \Cref{tab:instances} collects them.

\begin{table}[ht]
\centering
\small
\begin{tabular}{@{}p{0.17\linewidth} p{0.29\linewidth} p{0.23\linewidth} p{0.21\linewidth}@{}}
\toprule
\textbf{Domain} & \textbf{C2 violation (tilt $h$)} & \textbf{Margin $\gamma$} & \textbf{Singleton} \\
\midrule
Reliability \citep{towell2026mdrelax} &
  Group-structured informative masking &
  Masking-preference bound &
  A resolving diagnostic \\
scRNA-seq \citep{towell2026scrnacoarsening} &
  Dropout increasing in true expression &
  ERCC-vs-endogenous capture ratio &
  ERCC spike-in \\
Spatial \citep{towell2026spatialcoarsening} &
  Capture efficiency varying by cell type &
  Single-cell probe resolution &
  Single-cell-resolution probe \\
Differential privacy \citep{towell2026dpcoarsening} &
  Data-dependent mechanism (SVT, PTR) &
  Released kernel scale &
  Non-private release \\
Weak supervision \citep{towell2026weaksupcoarsening} &
  Labeling-function dependence &
  Accuracy gap &
  Gold-labeled example \\
Phenotyping \citep{towell2026phenotypecoarsening} &
  Severity-correlated coding &
  Case-mix gap &
  Chart-reviewed patient \\
\bottomrule
\end{tabular}
\caption{The C2-violation tilt, the identification margin, and the singleton in
each domain. \Cref{thm:sensitivity} reads the bias bound $\delta B(\theta)$ in the
tilt column; \cref{thm:restoration} reads the $r/\gamma^2$ rate in the margin
column. Reliability is the worked instance \citep{towell2026mdrelax}; the others
specialize the same two theorems.}
\label{tab:instances}
\end{table}

Two domains need a word. Differential privacy sits in the location-family regime of
\citet{towell2026synthesis}: its coarsening is engineered, so the tilt $h$ is not a
modeling unknown but a property of the chosen mechanism, and the C2 partition is
sharp, data-independent mechanisms (Laplace, Gaussian) have $\delta = 0$ by
construction, while data-dependent mechanisms (sparse vector, propose-test-release)
have $\delta > 0$ with $h$ determined by the realized query. The sensitivity bound
is therefore exact and computable rather than worst-case there, and the non-private
release is a singleton that both sets $\gamma$ and zeroes $\delta$. Weak supervision
is the cleanest illustration of \cref{thm:restoration}: the rank deficit $r$ is the
labeling-function dependence dimension and $\gamma$ is the accuracy gap, so the
gold-set rate $\Theta(r/\mathrm{gap}^2)$ of \citet{towell2026weaksupcoarsening} is
\eqref{eq:samplecomplexity} verbatim.

Reliability is the fully worked instance and ships the software: the relaxed
candidate-set tilt of \citet{towell2026mdrelax} is \cref{def:tilt} with a
group-structured $h$, its first-order robustness interval is the bias bound
\eqref{eq:Bdelta}, and its \texttt{ri\_first\_order} routine computes the identified
interval of \cref{cor:partial} for a user-supplied $\delta$.

The six subsections below specialize \cref{thm:sensitivity,thm:restoration} one
domain at a time, naming the tilt, the named bias bound the sibling already proves,
the rank deficit and margin, and the singleton. None re-derive a result: each reuses
the sibling's own constant.

\subsection{Reliability: masked component causes}
\label{sec:inst-reliability}

The latent $\Y$ is the failed component of a series system and the report is a
candidate set of components that could have caused the failure
\citep{towell2026masked}. C2 is non-informative masking: which admissible component
is named does not depend on which one actually failed. The tilt $h$ of
\cref{def:tilt} is a group-structured masking preference, realized in
\citet{towell2026mdrelax} as a perturbed masking-probability matrix $P(\delta) = P_0
+ \delta\, D$ with direction $D$, so $h(r,y)$ is the within-candidate-set log-weight
that favors some components over others. The sensitivity bound \eqref{eq:Bdelta}
specializes to that paper's \emph{first-order robustness interval}: its index of
local sensitivity to nonignorability (the slope of the estimand's deviation in
$\delta$ at $\delta = 0$) is the per-direction reading of $\Info^{-1}\Cov(s,h)$, and
its \texttt{ri\_first\_order} routine returns $\mathrm{tolerance}/|\mathrm{ISNI}|$,
the largest $\delta$ keeping the estimand inside a stated tolerance, which is exactly
the $\delta B(\theta^\star)$ radius of \cref{cor:partial} solved for $\delta$. For
\cref{thm:restoration} the rank deficit $r$ is the dimension of the confounded
subspace of the augmented candidate-set matrix of \citet{towell2026masked} (its
rank-deficiency witness has a one-dimensional null space along the
$(1,-1,-1,1)$ direction), and the margin $\gamma$ is the masking-preference signal a
resolving observation contributes along that subspace. The singleton is a resolving
diagnostic that names a single component, the identity-row that lifts the augmented
rank.

\subsection{Single-cell RNA sequencing: dropout}
\label{sec:inst-scrna}

The latent $\Y$ is a cell's true transcript count and the report is the observed
(zero-inflated) count \citep{towell2026scrnacoarsening}. C2 fails when the dropout
probability depends on the true expression level, so the tilt $h$ is dropout
increasing in $\Y$, parametrized there by a logit-intercept shift $s$ between the
spike-in and endogenous dropout curves; $s$ is our $\delta$. The sensitivity bound
\eqref{eq:Bdelta} is that paper's spike-in oracle bias bound (its
\texttt{thm:bias-bound}): the per-gene bias is $J_{\mu\phi}\,\Delta\phi + O(s^2)$
with $J_{\mu\phi} = -I_{\mu\mu}^{-1} I_{\mu\phi}$, which is precisely
$\Info^{-1}\Cov(s,h)$ written in the gene-level information blocks, and its
``useful zone'' $|s| \lesssim 0.5$ is the interval on which $\delta B(\theta^\star)$
stays below the naive-estimator error. Here the confounded subspace is
one-dimensional ($r = 1$, the dropout-shape direction the coarsened counts cannot
separate from the mean), and the margin $\gamma$ is the ERCC-versus-endogenous
capture ratio, the signal a spike-in carries about that direction. The singleton is
the ERCC spike-in, a transcript whose true count is known by construction; the
$r = 1$, finite-singleton case of \cref{thm:restoration} is the spike-in-fit bias
\citet{towell2026scrnacoarsening} reports.

\subsection{Spatial transcriptomics: cell-type deconvolution}
\label{sec:inst-spatial}

The latent $\Y$ is the cell-type composition of a spot and the report is the spot's
pooled expression \citep{towell2026spatialcoarsening}. C2 fails when capture
efficiency varies by cell type, so the tilt $h$ reweights the within-spot
contribution toward some types; the bias bound \eqref{eq:Bdelta} is that paper's
\emph{marker-gene bias bound} (its \texttt{thm:bias-marker}), in which the bias is
controlled by the conditioning of the marker matrix $M$, and which recovers
Tangram's empirical ``a few hundred marker genes suffice'' rule as the regime where
$\delta B(\theta^\star)$ is small. For \cref{thm:restoration} the rank deficit $r$
is the column-rank deficit of the spot-by-cell-type composition matrix $P$ (the
full-column-rank condition of its \texttt{thm:rank} is the coarsening-rank condition
of \citet{towell2026synthesis}, so $r = K - \operatorname{rank} P$ confounded
type directions), and the margin $\gamma$ is the resolution of a single-cell probe
along a deficient type. The singleton is a single-cell-resolution probe
(MERFISH / seqFISH+): a spot known to contain one cell type, the standard-basis row
that restores column rank.

\subsection{Differential privacy: data-dependent mechanisms}
\label{sec:inst-dp}

The latent $\Y$ is a confidential statistic and the report is its privatized release
\citep{towell2026dpcoarsening}. This domain sits in the location-family regime of
\citet{towell2026synthesis}, and the coarsening is engineered rather than observed,
so the tilt $h$ is not a modeling unknown but a property of the chosen mechanism: a
data-independent mechanism (Laplace, Gaussian) has $\delta = 0$ by construction,
while a data-dependent mechanism (sparse vector, propose-test-release) has $\delta >
0$ with $h$ fixed by the realized query path. \Cref{thm:sensitivity} is therefore
exact and computable here, not a worst-case envelope: the bias direction
$\Info^{-1}\Cov(s,h)$ is read off the known mechanism. For \cref{thm:restoration} the
deficient subspace is the one the composed Fisher informations fail to span jointly
(the span condition of its \texttt{thm:composition}, $r$ directions left unidentified
by an adaptive composition such as SVT), and the margin $\gamma$ is the released
kernel scale. The singleton is a non-private release of the statistic, which both
sets $\gamma$ and zeroes $\delta$ for the released coordinate.

\subsection{Weak supervision: labeling-function dependence}
\label{sec:inst-weaksup}

The latent $\Y$ is the true label and the report is the vector of labeling-function
(LF) votes \citep{towell2026weaksupcoarsening}. C2 fails when LFs are correlated
given the label, so the tilt $h$ is the LF-dependence structure. This is the
cleanest instance of \cref{thm:restoration}: the rank deficit $r$ is the
LF-dependence dimension, the margin $\gamma$ is the LF accuracy gap, and the
gold-set sample complexity $\Theta(r/\mathrm{gap}^2)$ proved as that paper's $T4$
(for total $L_2$ label-model recovery) is \eqref{eq:samplecomplexity} verbatim, down
to the loss-dependent refinement ($L_\infty$ recovery costs $\log r$ rather than
$r$). The singleton is a gold-labeled example, an instance whose true label is known,
the diagonal measurement that prices each confounded direction. On the
\cref{thm:sensitivity} side, the gold-free glass ceiling of
\citet{towell2026weaksupcoarsening} is the $r > 0$, $\delta$-unidentified endpoint
that the singleton budget collapses.

\subsection{Electronic phenotyping: informative coding}
\label{sec:inst-phenotype}

The latent $\Y$ is a patient's true clinical state and the report is the set of
diagnosis codes \citep{towell2026phenotypecoarsening}. C2 fails under informative
coding, when the probability of coding a true case increases with disease severity,
so the tilt $h$ is severity-correlated coding with scalar magnitude the coding-slope
$|b_1|$ (its $\delta$). The sensitivity bound \eqref{eq:Bdelta} is that paper's
\emph{informative-coding bias bound} (its \texttt{thm:bias-informative}): the gap
between the average-severity coding probability and the cohort sensitivity is at most
$|b_1|\,\sigma_{S,1}/4$, the explicit $\delta B(\theta^\star)$ in which the severity
spread $\sigma_{S,1}$ plays the role of the score-tilt covariance. For
\cref{thm:restoration} the confounded direction is the prevalence-sensitivity
trade-off the codes cannot resolve, and the margin $\gamma$ is the case-mix gap
between the chart-reviewed subsample and the cohort. The singleton is a
chart-reviewed patient, a record whose true state is adjudicated; the
chart-review-calibrated estimator of \citet{towell2026phenotypecoarsening} is the
Rogan-Gladen correction, the $r = 1$ instance of the singleton budget.

\section{Validation}
\label{sec:validation}

Two predictions are checked by simulation: that the bias bound \eqref{eq:Bdelta}
is tight (the realized bias grows linearly in $\delta$ at the predicted slope), and
that the singleton sample complexity follows the $r/\gamma^2$ rate of
\eqref{eq:samplecomplexity}. Both reuse the per-domain data-generating processes of
the sibling papers, instrumented with the single tilt knob $\delta$ of
\cref{def:tilt}; the shared harness is \texttt{scripts/sensitivity\_sweep.R}.

\paragraph{Sensitivity sweep.} For each domain we draw data under
$\mathrm{C2}(\delta)$ for a grid of $\delta$, fit the face-value MLE, and record the
realized parameter bias against the first-order prediction $\delta\,
\Info^{-1}\Cov(s, h)$ of \eqref{eq:bias}. The prediction is confirmed if the
realized bias is linear in $\delta$ through the origin with slope matching the
analytic constant, and if the realized identified set matches the ellipsoid of
\cref{cor:partial}. The exponential-family domains (single-cell, spatial,
phenotyping, the discrete weak-supervision parametrization) should track the
prediction with no $O(\delta^2)$ curvature along the curved directions; the
location-family domain (differential privacy) should track it up to the
$O_p(n^{-1/2})$ remainder.

\paragraph{Singleton sample complexity.} For a domain with a tunable rank deficit
$r$ and margin $\gamma$ we add $n_{\mathrm{s}}$ singletons and record the restored
parameter's squared error as a function of $n_{\mathrm{s}}$, $r$, and $\gamma$. The
prediction \eqref{eq:samplecomplexity} is confirmed if the error scales as
$1/n_{\mathrm{s}}$, the singletons needed for a fixed accuracy scale linearly in
$r$, and the dependence on the margin is $\gamma^{-2}$. The weak-supervision
$r$-sweep of \citet{towell2026weaksupcoarsening} is one already-reproduced slice of
this surface.

This section reports the harness output (figures and the fitted slopes) once the
sweep has been run across the six domains.

\section{Discussion}
\label{sec:discussion}

\paragraph{Two lineages the coarsening program now joins.} The sensitivity bound is
a structured instance of missing-not-at-random sensitivity analysis. The tilt
$\delta$ of \cref{def:tilt} is the coarsened-data counterpart of the selection tilt
of \citet{copas1997inference}, the identified set of \cref{cor:partial} is a Manski
bound \citep{manski1989anatomy}, and the impossibility of testing $\delta$ from the
observed data is the coarsened-data form of the equal-fit phenomenon of
\citet{molenberghs2008every}: two tilts with the same face-value fit and different
latent conclusions cannot be told apart without external information. What the
coarsening view adds to that literature is structure, the candidate set and the
coarsening-sufficient statistic, which turn an abstract sensitivity parameter into a
computable bias direction $\Info^{-1}\Cov(s, h)$ and an explicit constant
$B(\theta)$. The restoration theorem is the corresponding instance of
measurement-error double sampling and verification-bias correction
\citep{tenenbein1970double,begg1983assessment,carroll2006measurement}: the singleton
is internal validation, and \cref{thm:restoration} is its sample-complexity
accounting, with the margin $\gamma$ the analogue of validation-sample
informativeness.

\paragraph{Why C2 deserved this.} A coarsening-literate reader objects that C1, C2,
C3 are untestable and fragile \citep{grunwald2003updating,jaeger2005ignorability},
and that assuming them is the weak point of the program. This paper is the response:
it does not assume C2, it measures the cost of its failure and prices the cure. The
honest deliverable under informative coarsening is the identified set of
\cref{cor:partial} together with the singleton budget of \cref{thm:restoration},
not a point estimate presented as if C2 held.

\paragraph{Relation to the synthesis and to known identities.} With
\citet{towell2026synthesis} this completes a two-pillar account: that paper gives
the C2-holds identity, identifiability, and singleton cure; this one gives the
C2-fails bias bound, partial identification, and singleton sample complexity. The
exponential-family regularity that makes \cref{thm:sensitivity} exact along curved
directions is the coarsened-data moment-matching of \citet{domke2020moment} read
under a tilt; the rank deficit $r$ is the finite-support completeness gap of
\citet{neweypowell2003instrumental} and the single-view specialization of the
proxy-identification logic of \citet{miao2018identifying}; the singleton is the
substitute for the multiple conditionally independent views that the latent-variable
canon assumes and the coarsened-data setting lacks.

\paragraph{Limitations.} The bias bound is first order in $\delta$; the
exponential-family exactness covers the curved directions but the general
finite-$\delta$, finite-$n$ remainder is bounded, not characterized, and tightening
it (especially in the location-family regime) is open. The sample-complexity rate is
stated for a balanced singleton allocation over the confounded subspace; optimal
allocation when the margins differ across directions, and adaptive singleton
acquisition, are natural extensions. Estimating the rank deficit $r$ and the margin
$\gamma$ from data, rather than treating them as known, is the same open problem the
weak-supervision sibling records.

\section{Conclusion}
\label{sec:conclusion}

Coarsening at random is the exception; informative coarsening is the rule. This
paper prices the rule. When the symmetry condition C2 fails by a tilt of magnitude
$\delta$, the face-value maximum-likelihood fit is biased by an amount linear in
$\delta$ along a direction the marginal-fit check cannot see, so the latent
parameter is partially identified over a set of size order $\delta$ rather than a
point. A
singleton, the report that pins the latent value and is therefore classical
internal validation, restores point identification, and the number needed to
recover the $r$ confounded directions is of order $r/\gamma^2$ in the domain's
identification margin $\gamma$. The two results give the coarsening program its
missing half: the synthesis says what is recoverable when coarsening is ignorable,
and this paper says how recovery degrades when it is not, and exactly how much
external information buys it back. The reliability sensitivity bands, the
single-cell spike-in bias, the weak-supervision gold-set complexity, and the
differential-privacy mechanism partition are one sensitivity-and-restoration theory,
positioned in the lineage of MNAR sensitivity analysis and measurement-error double
sampling.


\bibliographystyle{imsart-nameyear}
\bibliography{refs}

\end{document}
